% paper/sections/10_full_algorithm.tex
% §10: 完全アルゴリズムフロー（リストラクチャ版）

\section{完全アルゴリズムフロー：離散形式}
\label{sec:algorithm}

第~\ref{sec:CCD}章から第~\ref{sec:ppe_solve}章まで，CCD 空間離散化，
CLS 界面追跡，時間積分，圧力連成，および Poisson ソルバーを個別に構築した．
本章ではこれらの構成要素を統合し，
時刻 $t^n$ から $t^{n+1}$ への\textbf{1タイムステップの完全アルゴリズム}を
離散形式でまとめる．

\begin{enumerate}
  \item \textbf{§~\ref{sec:algo_operator_map}}：NS 方程式の各項と離散ソルバーの対応を概観する．
  \item \textbf{§~\ref{sec:algo_flow}}：7 ステップの完全アルゴリズムを離散形式で詳述する．
  \item \textbf{§~\ref{sec:dccd_params}}：DCCD フィルタの散逸パラメータ設計と適応制御を与える．
  \item \textbf{§~\ref{sec:bootstrap}}：シミュレーション開始時の初期化シーケンスを示す．
  \item \textbf{§~\ref{sec:algo_timestep}}：タイムステップ制御の安定性制約を定める．
\end{enumerate}

% ═════════════════════════════════════════════════════════════════════════════
\subsection{演算子マッピングと全体概観}
\label{sec:algo_operator_map}
% ═════════════════════════════════════════════════════════════════════════════

図~\ref{fig:ns_solvers} に，NS 方程式の各項と本手法で用いる離散ソルバーの対応を示す．
コロケート格子（セル中心）上で圧力・速度・界面変数すべてが同一点に配置され，
DCCD 散逸によりチェッカーボード不安定を内在的に抑制する．

%%% ── 図：NS各項とソルバーの対応 ─────────────────────────────────────────
\begin{figure}[htbp]
\centering
\begin{tikzpicture}[
  node distance=4mm and 5mm,
  nsterm/.style={draw=navy!70, fill=lightblue, rounded corners=4pt,
                 font=\small\bfseries, inner sep=6pt,
                 minimum width=26mm, minimum height=10mm, align=center},
  solver/.style={draw=#1!70, fill=#1!10, rounded corners=4pt,
                 font=\scriptsize, inner sep=5pt,
                 minimum width=26mm, align=center},
  solver/.default=navy,
  subarr/.style={-Stealth, semithick, darkgray!60},
]

\node[font=\small\bfseries, text=navy] (hd1) at (0,0)
  {NS 方程式の各項とソルバーの対応};

%% ── NS 5項ボックス ─────────────────────────────────────────────────
\node[nsterm, below=7mm of hd1, xshift=-57mm] (T1) {対流項};
\node[nsterm, right=4mm of T1]  (T2) {圧力項};
\node[nsterm, right=4mm of T2]  (T3) {粘性項};
\node[nsterm, right=4mm of T3]  (T4) {重力項};
\node[nsterm, right=4mm of T4]  (T5) {表面張力項};

%% ── 各ソルバーボックス ─────────────────────────────────────────────
\node[solver=orange2, below=5mm of T1] (S1)
  {CCD（$D^{(1)}$）\\AB2（$\Ord{\Delta t^2}$）\\陽的};
\node[solver=midblue, below=5mm of T2] (S2)
  {欠陥補正法 PPE\\（DC，k=1）\\+ CCD 勾配\\Projection法};
\node[solver=teal2, below=5mm of T3] (S3)
  {CCD（$\Ord{h^6}$）\\Crank-Nicolson\\半陰的};
\node[solver=darkgray, below=5mm of T4] (S4)
  {陽的\\(定数項)};
\node[solver=red2, below=5mm of T5] (S5)
  {Level Set $\to\kappa$\\CCD\\陽的};

\foreach \t/\s in {T1/S1, T2/S2, T3/S3, T4/S4, T5/S5}
  \draw[subarr] (\t) -- (\s);

\end{tikzpicture}
\caption{NS 方程式の各項と対応ソルバーの概要．
対流項は CCD + AB2（陽的），粘性項は CCD + Crank--Nicolson（半陰的），
圧力項は FD 直接法 PPE + CCD 勾配（Projection），
表面張力は CSF 体積力として Predictor に含む．}
\label{fig:ns_solvers}
\end{figure}

DCCD フィルタの散逸パラメータ $\varepsilon_d$ は用途ごとに異なる値を取る
（設計根拠は §~\ref{sec:dccd_params} で詳述）：
\begin{itemize}[noitemsep]
  \item \textbf{CLS 移流}（§~\ref{sec:advection_motivation}）：$\varepsilon_d = 0.05$（界面プロファイル保持のため低散逸）
  \item \textbf{PPE 右辺・Corrector 発散}（§~\ref{sec:dccd_decoupling}）：$\varepsilon_d = 1/4$（チェッカーボード完全抑制）
  \item \textbf{曲率計算の DCCD フィルタ}（§~\ref{sec:curvature}）：$\varepsilon_d = 0.05$（高次微分の精度保持）
\end{itemize}

表~\ref{tab:algo_overview} に 7 ステップの概観を示す．

\begin{table}[htbp]
\centering
\caption{1 タイムステップ（$t^n \to t^{n+1}$）のアルゴリズム概観}
\label{tab:algo_overview}
\small
\begin{tabular}{clll}
\hline
Step & 処理 & 手法 & 参照 \\
\hline
1 & CLS 移流       & TVD-RK3 + DCCD（$\varepsilon_d=0.05$）       & §~\ref{sec:advection_motivation} \\
2 & 再初期化        & 仮想時間ステッピング                            & §~\ref{sec:cls_compression} \\
3 & 物性更新        & $\psi\!\to\!\phi$ logit + 線形補間              & §~\ref{sec:why_cls} \\
4 & 曲率計算        & CCD 高次微分                                   & §~\ref{sec:curvature} \\
5 & Predictor      & AB2 + CN/ADI + IPC + CSF                       & §~\ref{sec:ab2_ipc} \\
6 & PPE 求解       & DC + LU（$k\!=\!1$）                            & §~\ref{sec:ppe_solve} \\
7 & Corrector      & IPC 増分更新                                    & §~\ref{sec:pressure} \\
\hline
\end{tabular}
\end{table}

\begin{tcolorbox}[defbox, title={本章で用いる離散演算子の記法}]
\label{box:operator_notation}
\begin{itemize}[noitemsep]
  \item $D_x^{(1)}, D_y^{(1)}$：CCD 法による1次微分 $(\partial/\partial x,\, \partial/\partial y)$
  \item $D_x^{(2)}, D_y^{(2)}$：CCD 法による2次微分 $(\partial^2/\partial x^2,\, \partial^2/\partial y^2)$
  \item $D_{xy}^{(11)}$：CCD 法による混合微分 $\partial^2/\partial x\partial y$—
        $(D_{xy}^{(11)} f)_{i,j} = (D_y^{(1)}[D_x^{(1)} f])_{i,j}$（逐次2段スウィープ，$\Ord{h^6}$）
  \item $\mathcal{C}_{\text{CCD}}(\bu,\, q)$：CCD 中心差分移流オペレータ（NS 対流項専用）
          $= u\,D_x^{(1)} q + v\,D_y^{(1)} q$
          （非保存形；$\bnabla\cdot\bm{u}=0$ のもとで保存形と等価．
          CLS 移流には DCCD 保存形 $\mathrm{Div}_h^{\psi}$ を使用；§~\ref{sec:advection_motivation}）
  \item $\mathcal{L}_{\text{visc},x}$，$\mathcal{L}_{\text{visc},y}$：粘性応力テンソル $\bnabla\cdot\boldsymbol{\tau}$ の各成分（§~\ref{sec:time_int}）
  \item $\mathrm{Div}_h^{\psi}(\psi\bm{u})$：DCCD 保存形フラックス発散（§~\ref{sec:advection_motivation}）
  \item $\mathcal{L}_{\mathrm{CCD}}^{\rho}$：CCD product-rule 変密度ラプラシアン
          $= (1/\rho)(D_x^{(2)}p + D_y^{(2)}p)
            - (D_x^{(1)}\rho/\rho^2)D_x^{(1)}p
            - (D_y^{(1)}\rho/\rho^2)D_y^{(1)}p$
          （$\Ord{h^6}$；現行 k=1 では $\Ord{h^2}$，§~\ref{sec:ppe_discretization_choice}）
\end{itemize}
\end{tcolorbox}

% ═════════════════════════════════════════════════════════════════════════════
\subsection{完全アルゴリズム：離散形式}
\label{sec:algo_flow}
% ═════════════════════════════════════════════════════════════════════════════

時刻 $t^n$ から $t^{n+1} = t^n + \Delta t$ への計算手順を，離散化された形式で以下にまとめる．
添字 $(i,j)$ は格子点，上付き文字 $n, *, n+1$ は時間ステップを表す．

\begin{tcolorbox}[algbox, title={1 タイムステップ $t^n \to t^{n+1}$ の完全アルゴリズム}]
\label{alg:full_timestep}

% ── Step 1 ──────────────────────────────────────────────────────────────
\paragraph{Step 1：CLS 移流}
保存形レベルセット変数 $\psi^n$ を，速度場 $\bm{u}^n$ を用いて TVD-RK3 で時間発展させる：
\begin{align*}
  \psi^{(1)}_{i,j} &= \psi^n_{i,j} - \Delta t \, \mathrm{Div}_h^{\psi}(\psi^n \bm{u}^n)_{i,j} \\
  \psi^{(2)}_{i,j} &= \tfrac{3}{4}\psi^n_{i,j} + \tfrac{1}{4}\psi^{(1)}_{i,j} - \tfrac{1}{4}\Delta t \, \mathrm{Div}_h^{\psi}(\psi^{(1)} \bm{u}^n)_{i,j} \\
  \psi^*_{i,j} &= \tfrac{1}{3}\psi^n_{i,j} + \tfrac{2}{3}\psi^{(2)}_{i,j} - \tfrac{2}{3}\Delta t \, \mathrm{Div}_h^{\psi}(\psi^{(2)} \bm{u}^n)_{i,j}
\end{align*}
各ステージ完了後に $\psi \leftarrow \max(0,\min(1,\psi))$ のクランプを適用
（DCCD は TVD 性を持たないため必須；注意~\ref{warn:adv_clamp}）．

% ── Step 2 ──────────────────────────────────────────────────────────────
\paragraph{Step 2：再初期化}
移流後の $\psi^*$ を式~\eqref{eq:cls_reinit_iso} で整形する：
$\psi^{n+1} = \mathrm{Reinitialize}(\psi^*,\; n_{\mathrm{reinit}})$．
\begin{equation}
  \Delta\tau = \min\!\left(
    \underbrace{\frac{0.5\,h^2}{2\,N_{\mathrm{dim}}\,\varepsilon}}_{\text{放物安定条件}},\;
    \underbrace{0.5\,h}_{\text{双曲 CFL}}
  \right)
  \tag{\ref{eq:dtau_reinit_def} 再掲}
\end{equation}
デフォルト $n_{\mathrm{reinit}} = 4$（§~\ref{sec:cls_compression}，アルゴリズム~\ref{alg:cls_reinit_dccd}）．

% ── Step 3 ──────────────────────────────────────────────────────────────
\paragraph{Step 3：グリッド・物性更新}
$\psi^{n+1}$ に基づき距離関数と物性値を再計算する
（物性補間は $\psi$，幾何量は $\phi$；§~\ref{sec:notation}）：
\begin{align*}
  \phi^{n+1}_{i,j} &= \varepsilon\ln\!\left(\frac{\psi^{n+1}_{i,j}}{1-\psi^{n+1}_{i,j}}\right)
                    \quad (\text{ロジット関数，§~\ref{sec:newton_inversion}}) \\
  \rho^{n+1}_{i,j} &= \rho_l + (\rho_g-\rho_l)\,\psi^{n+1}_{i,j}, \qquad
  \mu^{n+1}_{i,j}  = \mu_l + (\mu_g-\mu_l)\,\psi^{n+1}_{i,j}
\end{align*}
\textbf{注：}$\rho^{n+1}$，$\mu^{n+1}$ は Step 4--7 で終始同一の値を使用する（中間更新なし）．

% ── Step 4 ──────────────────────────────────────────────────────────────
\paragraph{Step 4：曲率計算}
$\phi^{n+1}$ の CCD 微分から曲率を求める（式~\eqref{eq:curvature_2d}）：
\begin{align*}
 (\phi_x, \phi_y) &= (D_x^{(1)}\phi^{n+1},\; D_y^{(1)}\phi^{n+1}), \quad
 (\phi_{xx}, \phi_{yy}, \phi_{xy}) = (D_x^{(2)}\phi^{n+1},\; D_y^{(2)}\phi^{n+1},\; D_{xy}^{(11)}\phi^{n+1})
\end{align*}
\begin{equation*}
 \kappa^{n+1}_{i,j} = -\frac{\phi_y^2 \phi_{xx}
                              - 2\phi_x\phi_y\phi_{xy}
                              + \phi_x^2 \phi_{yy}}
                           {\max\!\bigl(\phi_x^2 + \phi_y^2,\;
                            \varepsilon_{\text{norm}}^2\bigr)^{3/2}}
\end{equation*}

% ── Step 5 ──────────────────────────────────────────────────────────────
\paragraph{Step 5：運動量予測（Predictor，AB2 + CN/ADI + IPC）}
対流項に AB2 外挿，粘性項に CN，IPC により $-\bnabla p^n$ を陽的に加え，
CSF 体積力 $(\kappa/We)\bnabla\psi$ を Predictor に含む
（Balanced-Force 条件，式~\eqref{eq:balanced_pressure}，\eqref{eq:balanced_surface}）．

\medskip
\underline{$x$ 成分（$u^*_{i,j}$，$n \ge 1$）：}
\begin{align*}
  \frac{\rho^{n+1}_{i,j}}{\Delta t}\bigl(u^*_{i,j} - u^n_{i,j}\bigr)
  &= -\rho^{n+1}_{i,j}\Bigl[
      \tfrac{3}{2}\mathcal{C}_{\text{CCD}}(\bu^n, u^n)
      -\tfrac{1}{2}\mathcal{C}_{\text{CCD}}(\bu^{n-1}, u^{n-1})
    \Bigr]_{i,j} \\
  &\quad + \frac{1}{2}\bigl[\mathcal{L}_{\text{visc},x}(\mu^{n+1}, \bu^*) + \mathcal{L}_{\text{visc},x}(\mu^{n+1}, \bu^n)\bigr]_{i,j} \\
  &\quad - \bigl(D_x^{(1)} p^n\bigr)_{i,j}
         + \frac{1}{We}\,\kappa^{n+1}_{i,j}\,\bigl(D_x^{(1)}\psi^{n+1}\bigr)_{i,j}
\end{align*}
$y$ 成分は $x \leftrightarrow y$ の対称形（付録~\ref{app:predictor_y}）．
スタートアップ（$n=0$）のみ対流項を前進 Euler で代用（§~\ref{sec:ab2_ipc}，注意~\ref{warn:ab2_startup}）．
ADI 分裂の実装詳細（陰的係数の選択，スウィープ形式）は付録~\ref{app:cn_adi_detail} 参照．

% ── Step 6 ──────────────────────────────────────────────────────────────
\paragraph{Step 6：圧力 Poisson 方程式の求解（IPC 増分法）}
圧力増分 $\delta p \equiv p^{n+1} - p^n$ を欠陥補正法で求解する：
\begin{equation}
  \left[\frac{D_x^{(2)}(\delta p)}{\rho^{n+1}}
  - \frac{D_x^{(1)}\rho^{n+1}}{(\rho^{n+1})^2}D_x^{(1)}(\delta p)\right]_{i,j}
  + \left[\frac{D_y^{(2)}(\delta p)}{\rho^{n+1}}
  - \frac{D_y^{(1)}\rho^{n+1}}{(\rho^{n+1})^2}D_y^{(1)}(\delta p)\right]_{i,j}
  = \frac{1}{\Delta t}\Bigl(D_x^{(1)}\tilde{u}^*_{i,j} + D_y^{(1)}\tilde{v}^*_{i,j}\Bigr)
  \label{eq:PPE_CCD_fullalgo}
\end{equation}
右辺は DCCD フィルタ後速度（$\varepsilon_d=1/4$，式~\eqref{eq:dccd_filter_physical}）の
CCD 発散から構成される（式~\eqref{eq:dccd_ppe_rhs}）．
境界条件：$\partial(\delta p)/\partial n = 0$（ノイマン条件）．
反復アルゴリズム（Steps 6.0--6.5）の詳細は付録~\ref{app:dc_ppe_detail} 参照．

\noindent\textbf{高密度比への拡張：}
上記の変密度 PPE（smoothed Heaviside 一括解法）は $\rho_l/\rho_g \leq 5$ で有効であるが，
$\rho_l/\rho_g \geq 10$ では DC が発散する（§~\ref{sec:interface_crossing}）．
高密度比では Step~6 を\textbf{分相 PPE 解法}——各相で定数密度 PPE を独立に解き，
界面ジャンプ条件を HFE（§~\ref{sec:field_extension}）で課す方式——に置換する．
§~\ref{sec:split_ppe_recovery} にて，分相 PPE が全密度比（$\rho_l/\rho_g = 1$--$1000$）で
密度比非依存な $\Ord{h^2}$ 収束を達成することを実証した．

% ── Step 7 ──────────────────────────────────────────────────────────────
\paragraph{Step 7：速度・圧力補正（Corrector，IPC）}
圧力増分 $\delta p$ を用いて速度と圧力を同時に更新する：
\begin{align*}
u^{n+1}_{i,j} &= u^*_{i,j}
  - \frac{\Delta t}{\rho^{n+1}_{i,j}}\,\bigl(D_x^{(1)} \delta p\bigr)_{i,j}, \quad
v^{n+1}_{i,j} = v^*_{i,j}
  - \frac{\Delta t}{\rho^{n+1}_{i,j}}\,\bigl(D_y^{(1)} \delta p\bigr)_{i,j} \\
p^{n+1}_{i,j} &= p^n_{i,j} + \delta p_{i,j}
\end{align*}
$\bm{u}^{n+1}$ は CCD 発散の意味で非圧縮条件を満たす：
$D_x^{(1)} u^{n+1} + D_y^{(1)} v^{n+1} \approx 0$
（DCCD フィルタが $2\Delta x$ モードを除去；§~\ref{sec:dccd_decoupling}）．
\end{tcolorbox}

% ═════════════════════════════════════════════════════════════════════════════
\subsection{DCCD パラメタの設計と適応制御}
\label{sec:dccd_params}
% ═════════════════════════════════════════════════════════════════════════════

§~\ref{sec:dccd_filter_theory} で導入した DCCD フィルタの理論的枠組みに基づき，
本節では散逸強度 $\varepsilon_{d,\max}$ の具体的な設計値と，
CLS 変数 $\psi$ に基づく適応制御アルゴリズムを与える．

表~\ref{tab:dccd_usage_map} に，アルゴリズム全体で DCCD フィルタが担う 3 つの役割と
パラメータ設定を一覧する．

\begin{table}[htbp]
\centering
\caption{DCCD フィルタのパラメータ設計一覧}
\label{tab:dccd_usage_map}
\small
\begin{tabular}{llcl}
\hline
用途 & $\varepsilon_d$ & 適応制御 & 設計根拠 \\
\hline
CLS 移流（§~\ref{sec:advection_motivation}）
  & 0.05 & あり（$S(\psi)$） & 界面プロファイル保持 \\
PPE RHS / Corrector（§~\ref{sec:dccd_decoupling}）
  & $1/4$ & なし（固定） & チェッカーボード完全抑制 \\
曲率計算（§~\ref{sec:curvature}）
  & 0.05 & あり（$S(\psi)$） & 高次微分の精度保持 \\
\hline
\end{tabular}
\end{table}

% ─────────────────────────────────────────────────────────────────────────────
\subsubsection{散逸強度 \texorpdfstring{$\varepsilon_{d,\max}$}{ε\_d,max} のスペクトル設計}
\label{sec:dccd_eps_design}
% ─────────────────────────────────────────────────────────────────────────────

スイッチ関数が決まれば，次は散逸強度の上限 $\varepsilon_{d,\max}$ を設計する問題となる．
設計の指針は以下の 2 つの要求の\textbf{同時達成}である：
\begin{enumerate}[label=(\alph*)]
  \item \textbf{精度保存：} 低波数域（長波長の物理情報）は CCD の $\Ord{h^6}$ 精度を損なわない．
  \item \textbf{高周波抑制：} ナイキスト近傍（$\xi \gtrsim 2\pi/3$，波長 $\lesssim 3h$）を
        有意に減衰させる．
\end{enumerate}

\medskip
\noindent\textbf{定量的設計．}

フィルタ伝達関数 $H(\xi;\varepsilon_d) = 1 - 4\varepsilon_d\sin^2(\xi/2)$ において，
低波数（$\xi \leq \xi_L$）での精度要求を：
\begin{equation}
  |H(\xi_L;\varepsilon_d) - 1| = 4\varepsilon_d\sin^2\!\left(\frac{\xi_L}{2}\right) \leq \delta_\mathrm{acc}
  \label{eq:accuracy_req}
\end{equation}
高波数（$\xi = \pi$）での減衰要求を：
\begin{equation}
  H(\pi;\varepsilon_d) = 1 - 4\varepsilon_d \leq 1 - D_\mathrm{target}
  \quad\Leftrightarrow\quad
  \varepsilon_d \geq \frac{D_\mathrm{target}}{4}
  \label{eq:damping_req}
\end{equation}
と定式化する．典型的な二相流計算では，$\xi_L = \pi/3$（波長 $\geq 6h$ を精度保存），
$\delta_\mathrm{acc} = 0.05$（5\% 以内の振幅変化を許容），
$D_\mathrm{target} = 0.20$（ナイキスト成分を 20\% 減衰）とする：

\begin{tcolorbox}[resultbox, title={$\varepsilon_{d,\max}$ の設計値と設計根拠}]
\label{box:eps_design}
精度要求（式~\eqref{eq:accuracy_req}，$\xi_L = \pi/3$，$\delta_\mathrm{acc}=0.05$）から：
\begin{equation}
  4\varepsilon_d\sin^2(\pi/6) = 4\varepsilon_d\cdot\frac{1}{4} = \varepsilon_d \;\leq\; 0.05
  \label{eq:eps_bound_accuracy}
\end{equation}

減衰要求（式~\eqref{eq:damping_req}，$D_\mathrm{target}=0.20$）から：
\begin{equation}
  \varepsilon_d \;\geq\; \frac{0.20}{4} = 0.05
  \label{eq:eps_bound_damping}
\end{equation}

両制約が $\varepsilon_{d,\max} = 0.05$ で\textbf{等号達成}：

\begin{equation}
  \boxed{\varepsilon_{d,\max} = 0.05}
  \label{eq:eps_max}
\end{equation}

\textbf{解釈：}
波長 $6h$ 以上（物理的に意味ある成分）の精度劣化は最大 5\%，
ナイキスト成分（波長 $2h$）は 20\% 減衰する．
より強い安定化が必要な場合（高密度比・高 Reynolds 数計算）は
$\varepsilon_{d,\max} = 0.10$〜$0.20$ まで拡張できるが，
この場合は波長 $6h$ の精度劣化も 10〜20\% に増加することに注意する．

\medskip
\textbf{安定性制約の確認：}
$\varepsilon_{d,\max} = 0.05 \leq 0.25$（式~\eqref{eq:dccd_transfer} の安定限界 $1/4$）を満足する．$\square$
\end{tcolorbox}

\noindent\textbf{フィルタ伝達関数の波数プロファイル（$\varepsilon_d = 0.05$）：}
\begin{align*}
  H(0;\, 0.05)     &= 1.00 \quad (\xi = 0,\text{ DC 成分，変化なし}) \\
  H(\pi/3;\, 0.05) &= 1 - 4(0.05)\sin^2(\pi/6) = 1 - 0.05 = 0.95 \quad (\text{波長}\ 6h,\ 5\%\ \text{減衰}) \\
  H(2\pi/3;\, 0.05)&= 1 - 4(0.05)\sin^2(\pi/3) = 1 - 0.15 = 0.85 \quad (\text{波長}\ 3h,\ 15\%\ \text{減衰}) \\
  H(\pi;\, 0.05)   &= 1 - 4(0.05) = 0.80 \quad (\text{波長}\ 2h,\ 20\%\ \text{減衰})
\end{align*}

% ─────────────────────────────────────────────────────────────────────────────
\subsubsection{適応散逸制御則と実装アルゴリズム}
\label{sec:dccd_adaptive}
% ─────────────────────────────────────────────────────────────────────────────

\textbf{適応散逸強度．}
格子点 $i$ での散逸強度を CLS 変数 $\psi_i^n$（時刻 $t^n$）によって点別に制御する：
\begin{equation}
  \varepsilon_d^{(i)} \;=\; \varepsilon_{d,\max} \cdot S(\psi_i^n)
  \;=\; \varepsilon_{d,\max}\,(2\psi_i^n - 1)^2
  \label{eq:adaptive_eps}
\end{equation}
これにより：
\begin{itemize}
  \item $\psi_i \approx 0.5$（界面）：$\varepsilon_d^{(i)} \approx 0$，フィルタ無効，標準 CCD に退化
  \item $\psi_i \approx 0$ または $1$（バルク）：$\varepsilon_d^{(i)} \approx \varepsilon_{d,\max}$，
        最大散逸でノイズを積極的に排除
\end{itemize}

\begin{tcolorbox}[algbox, title={適応 Dissipative-CCD フィルタの実装アルゴリズム}]
\label{alg:dccd}
\textbf{入力：} スカラー場 $f$ の格子値 $\{f_i\}$，CLS 変数 $\{\psi_i^n\}$，散逸強度 $\varepsilon_{d,\max}$

\begin{enumerate}
  \item \textbf{CCD 求解：}
        標準 CCD ブロック Thomas アルゴリズム（§~\ref{sec:ccd_matrix}）を実行し，
        $\{f'_i\}$，$\{f''_i\}$ を得る．
  \item \textbf{点別散逸強度の計算：}
        各格子点 $i$ に対して
        \[
          \varepsilon_d^{(i)} = \varepsilon_{d,\max}\,(2\psi_i^n - 1)^2
        \]
        を計算する（計算コスト：$O(N)$，数値演算 1 回/点）．
  \item \textbf{1 次微分のフィルタ適用：}
        \[
          \tilde{f}'_i = f'_i + \varepsilon_d^{(i)}\,(f'_{i+1} - 2f'_i + f'_{i-1}),
          \quad i = 1, \ldots, N-2
        \]
        境界点 $i=0, N-1$：周期 BC では折り返しインデックスを使用；
        壁面・流出 BC では境界点をフィルタ適用前の CCD 値のまま保持する（§~\ref{sec:dccd_bc} 参照）．
  \item \textbf{2 次微分のフィルタ適用：}
        \[
          \tilde{f}''_i = f''_i + \varepsilon_d^{(i)}\,(f''_{i+1} - 2f''_i + f''_{i-1}),
          \quad i = 1, \ldots, N-2
        \]
  \item \textbf{後続計算への引き渡し：}
        フィルタ適用済みの $(\tilde{f}', \tilde{f}'')$ を曲率計算・速度勾配評価・
        圧力勾配計算（第~\ref{sec:pressure}章）に用いる．
\end{enumerate}

\textbf{2 次元への拡張：}
2 次元計算では $x$ 方向と $y$ 方向に独立に 1 次元フィルタを適用する（次元分割法）：
\[
  \tilde{f}'_{x,ij} = f'_{x,ij} + \varepsilon_d^{(ij)}\,(f'_{x,i+1,j} - 2f'_{x,ij} + f'_{x,i-1,j}),
  \quad
  \tilde{f}'_{y,ij} = f'_{y,ij} + \varepsilon_d^{(ij)}\,(f'_{y,i,j+1} - 2f'_{y,ij} + f'_{y,i,j-1})
\]
$f''_{xx}, f''_{yy}, f''_{xy}$ に対しても同様に適用する．

\textbf{計算コスト：} ステップ 1（CCD，$O(N^2)$ for 2D）+ ステップ 2--4（フィルタ，$O(N^2)$）．
フィルタは CCD の後処理として追加コストが小さい．
また，$\tilde{f}'$ と $\tilde{f}''$ はフィルタ段階で\textbf{独立}に補正する——
すなわち $\tilde{f}''$ は $\tilde{f}'$ から再連立せずに計算する（計算コスト削減のための設計選択）．
\end{tcolorbox}

% ═════════════════════════════════════════════════════════════════════════════
\subsection{ブートストラップ処理：初期化シーケンス}
\label{sec:bootstrap}
% ═════════════════════════════════════════════════════════════════════════════

シミュレーション開始前に，以下の初期化シーケンス（ブートストラップ処理）を一度だけ実行する．
格子生成には Level Set 場 $\phi^{(0)}$ が必要だが，$\phi^{(0)}$ の評価には格子が必要という
循環参照を，一様格子による仮評価で解消する．

\begin{tcolorbox}[algbox, title={ブートストラップ処理：初期化シーケンス}]
\begin{enumerate}
  \item \textbf{一様格子上での $\phi^{(0)}$ 計算：}
        $N_x \times N_y$ 等間隔格子上で $\phi^{(0)}_{i,j}$ を計算する
        （界面形状が解析的に与えられる場合は直接代入；そうでない場合は
        一様格子上の符号付き距離関数として初期化）．
  \item \textbf{非一様格子の生成（§~\ref{sec:grid_2d}）：}
        $\phi^{(0)}$ から界面適合非一様格子を生成する：
        \begin{enumerate}
          \item $x$ 方向の密度関数 $\omega^x_i = 1 + (\alpha-1)\,\varepsilon_g\,\delta_\varepsilon^*(\bar{\phi}^x_i)$ を計算し，
                §~\ref{sec:grid_gen} のステップ2--5で物理座標 $\{x_i\}$ とメトリクス係数 $\{J_x|_i\}$ を得る．
          \item $y$ 方向も同様に $\{y_j\}$ と $\{J_y|_j\}$ を得る．
        \end{enumerate}
  \item \textbf{CLS 変数の初期化：}
        非一様格子上で $\psi^{(0)}_{i,j} = H_\varepsilon(\phi^{(0)}_{i,j})$ を評価し，
        必要に応じて CLS 再初期化（§~\ref{sec:reinit}）を数ステップ実行して
        $|\bnabla\phi| \approx 1$ を確保する．
  \item \textbf{CCD 演算子の構築：}
        非一様格子のメトリクス係数を用いて CCD ブロック三重対角行列
        （§~\ref{sec:ccd_matrix}）を組み立て，LU 分解を事前計算する．
  \item \textbf{初期圧力場の設定：}
        静的平衡条件 $p^{(0)} = \sigma\kappa_0$（Young--Laplace 圧力）を初期値として設定する．
        曲率 $\kappa_0$ は $\psi^{(0)}$ から CCD 微分を用いて評価する．
\end{enumerate}

\textbf{出力：}
非一様格子 $\{(x_i, y_j)\}$，メトリクス係数 $\{J_x|_i, J_y|_j\}$，
CLS 場 $\psi^{(0)}$，CCD 演算子（LU 分解済み），初期圧力場 $p^{(0)}$．
\end{tcolorbox}

格子の時間更新モード（Mode 1: 固定格子，Mode 2: 再生成）については
§~\ref{sec:grid_ale} を参照されたい．

% ═════════════════════════════════════════════════════════════════════════════
\subsection{タイムステップ制御}
\label{sec:algo_timestep}
% ═════════════════════════════════════════════════════════════════════════════

各タイムステップの開始時に，以下の2制約のうち最も厳しい条件で $\Delta t$ を決定する
（詳細な導出は §~\ref{sec:stability} 参照）：

\begin{equation}
  \Delta t = C_{\mathrm{CFL}} \cdot \min\!\left(
    \underbrace{\left[\max\!\left(\frac{|u|_{\max}}{\Delta x}+\frac{|v|_{\max}}{\Delta y}\right)\right]^{-1}}_{\Delta t_{\mathrm{adv}}\text{：対流 CFL（式~\eqref{eq:dt_adv}）}},\;
    \underbrace{\sqrt{\frac{(\rho_l+\rho_g)\,h_{\min}^3}{2\pi\sigma}}}_{\Delta t_\sigma\text{：毛管波制約（式~\eqref{eq:dt_sigma}）}}
  \right)
  \label{eq:dt_full_algo}
\end{equation}

ここで $C_{\mathrm{CFL}} \le 0.5$ は安全率（デフォルト 0.45），$h_{\min} = \min(\Delta x, \Delta y)$ である．
粘性項は Crank--Nicolson 法（無条件安定）で処理するため，等粘性系では粘性 CFL 制約は不要である．
ただし大密度比（$\mu_l/\mu_g \gg 1$）の気液系では，クロス微分項 $\partial_y(\mu\partial_x v)$ の
陽的処理に起因する実質的な粘性制約が界面近傍で復活する場合がある（§~\ref{sec:stability} 参照）．

\bigskip
以上の統合アルゴリズムの各構成要素が設計通りの精度を達成するか，
次章で系統的にコンポーネント単位の数学的検証を行う．
